\documentclass[12pt,a4paper]{article}
\usepackage{goyabase}

\usepackage{longtable}
\usepackage{calc}

% Pandoc compatibility
\providecommand{\tightlist}{%
  \setlength{\itemsep}{0pt}\setlength{\parskip}{0pt}}
\newcommand{\passthrough}[1]{#1}

\title{\textbf{Redis en GloboTour: Velocidad y Tiempo Real}}
\author{Departamento de Informática}
\date{\today}

\usepackage[spanish, es-noshorthands, es-tabla]{babel}
\usepackage[autostyle]{csquotes}
\begin{document}

\maketitle

\hypertarget{redis-en-globotour-velocidad-y-tiempo-real}{%
\section{Redis en GloboTour: Velocidad y Tiempo
Real}\label{redis-en-globotour-velocidad-y-tiempo-real}}

En \textbf{GloboTour}, usamos Redis para gestionar las partes de la
aplicación que no pueden esperar a que un disco duro gire. Es nuestra
capa de aceleración.

\begin{center}\rule{0.5\linewidth}{0.5pt}\end{center}

\hypertarget{nomenclatura-y-jerarquuxeda-el-uso-de-los}{%
\subsection{\texorpdfstring{Nomenclatura y Jerarquía (El uso de los
\texttt{:})}{1. Nomenclatura y Jerarquía (El uso de los :)}}\label{nomenclatura-y-jerarquuxeda-el-uso-de-los}}

Redis es una base de datos \textbf{plana} (Key-Value), lo que significa
que no tiene tablas ni carpetas reales. Sin embargo, usamos los dos
puntos \texttt{:} por convención para organizar
nuestras llaves.

\begin{itemize}
\tightlist
\item
  \textbf{¿Qué son los \texttt{:}?}: Son un simple
  separador de texto. Para Redis, la llave
  \texttt{session:user:101} es solo una cadena de
  texto larga.
\item
  \textbf{Simulación de Carpetas}: Herramientas como
  \textbf{RedisInsight} interpretan los \texttt{:} para
  agrupar llaves visualmente en carpetas, facilitando la administración.
\item
  \textbf{Buenas prácticas}: Se recomienda seguir un patrón
  \texttt{entidad:id:atributo}.

  \begin{itemize}
  \tightlist
  \item
    Ejemplo: \texttt{exp:500:nombre} o
    \texttt{user:101:session\_token}.
  \end{itemize}
\end{itemize}

\begin{center}\rule{0.5\linewidth}{0.5pt}\end{center}

\hypertarget{visualizaciuxf3n-con-redisinsight-gui}{%
\subsection{Visualización con RedisInsight
(GUI)}\label{visualizaciuxf3n-con-redisinsight-gui}}

Para gestionar la base de datos de forma visual, sigue estos pasos:

\begin{enumerate}
\def\labelenumi{\arabic{enumi}.}
\tightlist
\item
  Abre tu navegador en \url{http://localhost:5540}.
\item
  Si es la primera vez, haz clic en \textbf{``Add Redis Database''}.
\item
  \textbf{Configuración CRÍTICA}:

  \begin{itemize}
  \tightlist
  \item
    \textbf{Host}: \texttt{redis-server} (esto permite que el
    contenedor de la web hable con el del servidor). \textbf{NO} uses
    \texttt{localhost} ni
    \texttt{127.0.0.1} aquí.
  \item
    \textbf{Port}: \texttt{6379}
  \item
    \textbf{Database Alias}: ``BBDD\_Clase''
  \item
    \textbf{User/Password}: Déjalos \textbf{vacíos}.
  \end{itemize}
\item
  Una vez añadida, entra en la sección \textbf{``Browser''} del menú
  izquierdo.
\item
  \textbf{Verás las ``Carpetas''}: Gracias a los
  \texttt{puntos (:)}, verás los datos organizados por
  \texttt{session}, \texttt{exp},
  \texttt{user}, etc. Puedes hacer clic en ellas para
  expandirlas y ver las llaves.
\end{enumerate}

\begin{center}\rule{0.5\linewidth}{0.5pt}\end{center}

\hypertarget{compendio-de-tipos-de-datos-buxe1sicos}{%
\subsection{Compendio de Tipos de Datos
Básicos}\label{compendio-de-tipos-de-datos-buxe1sicos}}

\begin{longtable}[]{@{}
  >{\raggedright\arraybackslash}p{(\columnwidth - 4\tabcolsep) * \real{0.3333}}
  >{\raggedright\arraybackslash}p{(\columnwidth - 4\tabcolsep) * \real{0.3333}}
  >{\raggedright\arraybackslash}p{(\columnwidth - 4\tabcolsep) * \real{0.3333}}@{}}
\toprule\noalign{}
\begin{minipage}[b]{\linewidth}\raggedright
Tipo de Dato
\end{minipage} & \begin{minipage}[b]{\linewidth}\raggedright
Descripción
\end{minipage} & \begin{minipage}[b]{\linewidth}\raggedright
Caso de Uso Típico
\end{minipage} \\
\midrule\noalign{}
\endhead
\bottomrule\noalign{}
\endlastfoot
\textbf{Strings} & El más simple. Texto, números o binarios (hasta
512MB). & Sesiones, contadores, caché HTML. \\
\textbf{Hashes} & Mapas de campos-valores (objetos). & Perfiles de
usuario, productos. \\
\textbf{Lists} & Listas de strings ordenadas por inserción. & Colas de
tareas (LIFO/FIFO). \\
\textbf{Sets} & Colecciones de elementos únicos (sin orden). &
Etiquetas, amigos en común, votaciones. \\
\textbf{Sorted Sets} & Elementos únicos donde cada uno tiene una
puntuación (\emph{score}). & Rankings, tablas de clasificación. \\
\end{longtable}

\begin{center}\rule{0.5\linewidth}{0.5pt}\end{center}

\hypertarget{ejemplos-pruxe1cticos-por-tipo-de-dato}{%
\subsection{Ejemplos Prácticos por Tipo de
Dato}\label{ejemplos-pruxe1cticos-por-tipo-de-dato}}

\hypertarget{a.-gestiuxf3n-de-sesiones-strings-con-ttl}{%
\subsubsection{Gestión de Sesiones (Strings con
TTL)}\label{a.-gestiuxf3n-de-sesiones-strings-con-ttl}}

\textbf{Comandos:} \texttt{SET} (crear),
\texttt{GET} (leer), \texttt{EXPIRE}
(expiración).

\begin{lstlisting}
# Crear sesion que dura 30 segundos
SET sesion:token:9988 "user_id:105" EX 30

# Consultar valor
GET sesion:token:9988
\end{lstlisting}

\hypertarget{b.-contadores-atuxf3micos-strings}{%
\subsubsection{Contadores Atómicos
(Strings)}\label{b.-contadores-atuxf3micos-strings}}

\textbf{Comandos:} \texttt{INCR} (sumar 1),
\texttt{INCRBY} (sumar N).

\begin{lstlisting}
# Incrementar visitas de una pagina
INCR metricas:visitas:inicio
INCRBY metricas:visitas:oferta:paris 10
\end{lstlisting}

\hypertarget{c.-perfiles-de-usuario-hashes}{%
\subsubsection{Experiencias del Catálogo
(Hashes)}\label{c.-perfiles-de-usuario-hashes}}

\textbf{Comandos:} \texttt{HSET} (asignar campos),
\texttt{HGET} (leer uno),
\texttt{HGETALL} (leer todo).

\begin{lstlisting}
# Guardar o modificar una experiencia (ya existe la 500)
HSET exp:500 nombre "Vuelo en Globo Segovia" precio 150 stock 5

# Obtener solo el precio
HGET exp:500 precio

# Consultar toda la informacion de la experiencia
HGETALL exp:500

# Incrementar el stock (vuelven a tener plazas)
HINCRBY exp:500 stock 2
\end{lstlisting}

\hypertarget{d.-colas-de-trabajo-lists}{%
\subsubsection{Historial de Búsquedas
(Lists)}\label{d.-colas-de-trabajo-lists}}

\textbf{Comandos:} \texttt{LPUSH} (insertar al inicio),
\texttt{RPOP} (extraer del final),
\texttt{LRANGE} (ver rango).

\begin{lstlisting}
# Añadir busquedas recientes
LPUSH busquedas:recientes "Paris"
LPUSH busquedas:recientes "Roma"

# Ver las 3 ultimas busquedas
LRANGE busquedas:recientes 0 2

# Extraer la busqueda mas antigua para limpieza
RPOP busquedas:recientes
\end{lstlisting}

\hypertarget{e.-rankings-y-leaderboards-sorted-sets}{%
\subsubsection{Ranking de Ventas (Sorted
Sets)}\label{e.-rankings-y-leaderboards-sorted-sets}}

\textbf{Comandos:} \texttt{ZADD} (añadir con puntos),
\texttt{ZREVRANGE} (ver top),
\texttt{ZRANK} (ver posición).

\begin{lstlisting}
# Añadir o actualizar ventas (ya existen en datos_iniciales)
ZADD experiencias:top_ventas 150 "Vuelo en Globo Segovia"
ZADD experiencias:top_ventas 320 "Cata de Vinos Ribera"

# Obtener las 2 mas vendidas (con sus ventas/score)
ZREVRANGE experiencias:top_ventas 0 1 WITHSCORES
\end{lstlisting}

\hypertarget{f.-etiquetas-y-conjuntos-sets}{%
\subsubsection{Intereses de Usuarios
(Sets)}\label{f.-etiquetas-y-conjuntos-sets}}

\textbf{Comandos:} \texttt{SADD} (añadir único),
\texttt{SINTER} (intersección),
\texttt{SUNION} (unión).

\begin{lstlisting}
# Asignar intereses a usuarios
SADD user:105:intereses "aventura" "aire_libre" "fotografia"
SADD user:202:intereses "gastronomia" "relax" "aventura"

# ¿Que intereses tienen en comun para una oferta conjunta?
SINTER user:105:intereses user:202:intereses
\end{lstlisting}

\begin{center}\rule{0.5\linewidth}{0.5pt}\end{center}

\section{Desafíos de Aprendizaje (Para practicar en clase)}

\begin{enumerate}
    \item \textbf{El Tesoro Efímero}: Crea una llave llamada \texttt{secreto} con el valor \texttt{oro} que dure solo 15 segundos. Intenta leerla a los 5 segundos y a los 20. ¿Qué ocurre? (Uso de \texttt{TTL}).
    \item \textbf{Historial de Navegación}: Usa la lista \texttt{busquedas:recientes} para añadir dos nuevas ciudades. Asegúrate de que solo se vean las 3 más recientes usando \texttt{LRANGE}.
    \item \textbf{El Ranking Dinámico}: En el Sorted Set de \texttt{experiencias:top\_ventas}, incrementa las ventas de \enquote{Cata de Vinos Ribera} en 100 usando \texttt{ZINCRBY}. ¿Se ha movido de posición automáticamente en RedisInsight?
\end{enumerate}


\end{document}
